\documentclass[12pt]{article}
\usepackage[utf8]{inputenc}
\usepackage[T1]{fontenc}
\usepackage[margin=1in]{geometry}
\usepackage{setspace}
\usepackage{xcolor}
\usepackage[colorlinks=true, urlcolor=blue!60!black, linkcolor=blue!60!black]{hyperref}
\setlength{\parindent}{0em}
\setlength{\parskip}{0.6em}
\renewcommand{\baselinestretch}{1.0}
\begin{document}
\begin{center}
{\LARGE \textit{Moore v. City of East Cleveland}, 431 U.S. 494 (1977)\par}
\vspace{0.5em}
{\large Case Brief --- Two-Page Edition\par}
\end{center}

\section{Memory Jogger}
A grandmother's conviction for living with two grandsons who were cousins was reversed 5--4; the plurality protected extended-family living as deeply rooted in history and tradition (\href{https://supreme.justia.com/cases/federal/us/431/494/}{Justia}).

\section{Detailed Case Facts}
East Cleveland defined ``family'' to allow at most one dependent child's children in a household. Inez Moore lived with her son and two grandsons who were first cousins; one had come to her after his mother died. The city ordered her to remove the 10-year-old as an ``illegal occupant,'' and when she refused, prosecuted her.

\section{Procedural History}
She was convicted and sentenced to five days in jail and a \$25 fine. The Ohio Court of Appeals affirmed under \textit{Belle Terre}, and the Ohio Supreme Court dismissed her appeal.

\section{Judicial Votes}
Reversed, 5--4. Powell wrote for Brennan, Marshall, and Blackmun; Brennan (with Marshall) concurred; Stevens concurred in the judgment; Burger, Stewart (with Rehnquist), and White dissented.

\section{Holding}
The ordinance violates the Fourteenth Amendment. For the plurality, when government intrudes on family living arrangements, courts closely examine its interests and the fit of its means. Stevens rested on the owner's right to use her property free of arbitrary zoning.

\section{Analysis of Opinions}
\textbf{Plurality:} Unlike \textit{Belle Terre}, which limited only unrelated occupants, this law intruded on the family itself, and its relation to crowding, traffic, and school costs was ``tenuous.'' Limits on substantive due process come from history and basic values, not an arbitrary line at the nuclear family. \textbf{Brennan:} the nuclear-family model reflects ``cultural myopia.'' \textbf{Stevens:} no zoning precedent allowed barring an owner's own relatives. \textbf{Dissents:} Burger would require a variance application first; Stewart found no fundamental right; White warned against judge-made rights. Powell's test was later adopted in \textit{Glucksberg} and \textit{Dobbs}.

\section{Examples: Future Applications}
\textbf{Same-Side:} prosecuting a woman for housing her elderly aunt and widowed sister is unconstitutional. \textbf{Opposite-Side:} a cap of three unrelated roommates stands under \textit{Belle Terre}. \textbf{Fence-Sitter:} a long-term chosen family including an unadopted godchild.

\section{Critique}
Originalists say the history-and-tradition test leaves judges free to choose the level of generality, though they later embraced it as a limit. Progressives praise its respect for kinship households but fault its blood-based definition of family.

\section{Key Quotations}
\begin{quote}
``Our decisions establish that the Constitution protects the sanctity of the family precisely because the institution of the family is deeply rooted in this Nation's history and tradition.''
\end{quote}
\noindent --- Justice Powell, plurality opinion, 431 U.S. at 503.

\end{document}
